\chapter{Contexto: a lei e a Chapada do Apodi}
\label{sec:politica}

A Chapada do Apodi é um planalto contínuo que a divisa entre o Ceará e o Rio
Grande do Norte corta ao meio. O relevo, o regime de chuvas e o aquífero de
que dependem a irrigação e o abastecimento são os mesmos dos dois lados; muda a jurisdição. Do lado cearense, o Perímetro Irrigado
Jaguaribe--Apodi, em Limoeiro do Norte e Quixeré, concentrou a partir dos anos
2000 a fruticultura de exportação. Essa unidade física dividida por uma fronteira legal
é o experimento natural do projeto.

A Lei Estadual nº~16.820, de 8 de janeiro de 2019, acrescentou o art.~28-B à
lei cearense de agrotóxicos (Lei nº~12.228/1993) e vedou a pulverização aérea
de agrotóxicos na agricultura em todo o estado, a partir da publicação, em 9 de
janeiro. Em 19 de dezembro de 2024 a Lei nº~19.135 abriu exceção para a
aplicação por drones. O Rio
Grande do Norte permaneceu todo o tempo sob a regra federal, que não proíbe o
avião e fixa distâncias mínimas: quinhentos metros de povoações e de
mananciais de captação, duzentos e cinquenta metros de outros mananciais e de
moradias isoladas \cite{mapa2008in2}. Um projeto de proibição foi apresentado
à Assembleia potiguar em 2019 \cite{alrn2019pl}, e o Ceará segue como o único
estado com lei em vigor \cite{olivete2025}.

Na Chapada cearense o avião servia à banana, contra a sigatoka, doença fúngica
das folhas da bananeira cujo controle exige aplicações repetidas de fungicida
na quadra chuvosa. A calda era de cerca de 25 litros por hectare, e os
produtos, fungicidas como difenoconazol, propiconazol e tebuconazol
\cite[p.~526]{teixeira2011rebulico}. A literatura econômica mediu sobretudo
herbicidas em grãos; aqui o caso é de fungicida em fruticultura perene, como
em \citeonline{calzada2023bananas}. Sobre a escala há um indício: o trabalho de campo de 2009
registrou, por informação do técnico e planilha da empresa, um contrato de 1.000 hectares, dos quais 820 de duas
empresas e 180 de vinte pequenos produtores parceiros
\cite[p.~171]{marinho2011dimensao}.

O Censo Agropecuário de 2006 dá a linha de base por município
(Tabela~\ref{tab:censo}). Dos 36 estabelecimentos cearenses que declararam
aplicação por aeronave, 27 estavam em Limoeiro do Norte e Quixeré, a maioria
não familiares. O avião não aparece no lado potiguar da Chapada.

\begin{table}[htb]
\centering\small
\caption{Estabelecimentos que usaram agrotóxico, por equipamento de aplicação,
2006}
\label{tab:censo}
\begin{tabular}{@{}lrrrrr@{}}
\toprule
 & \multicolumn{3}{c}{Por aeronave} & Tração & \\
\cmidrule(lr){2-4}
Município & Total & Não familiar & Familiar & mecânica & Costal \\
\midrule
Limoeiro do Norte (CE) & 18 & 14 & 4 & 56 & 871 \\
Quixeré (CE)           &  9 &  8 & 1 & 15 & 506 \\
Jaguaruana (CE)        &  1 &  1 & 0 &  9 & 959 \\
Russas (CE)            &  0 &  0 & 0 & 18 & 2.020 \\
Apodi (RN)             &  0 &  0 & 0 & 30 & 1.503 \\
Baraúna (RN)           &  0 &  0 & 0 & 107 & 418 \\
Mossoró (RN)           &  0 &  0 & 0 & 15 & 546 \\
\midrule
Ceará, total           & 36 & 24 & 12 & & \\
Rio Grande do Norte, total & 5 & 0 & 5 & & \\
\bottomrule
\end{tabular}
\fonte{IBGE, Censo Agropecuário 2006, tabela SIDRA 2069, extraída em 16/11/2026. Um estabelecimento pode declarar mais de um equipamento. O Censo de 2017 não repetiu a pergunta.}
\end{table}

\section{O que não se sabe sobre os anos anteriores à lei}
\label{sec:quando}

Antes do desfecho, um desenho de diferenças em diferenças deve mostrar que a
política alterou o tratamento \cite{cunningham2021mixtape}, e aqui esse
primeiro estágio está por medir. A pulverização aérea na Chapada está
documentada com medida até 2010 (Figura~\ref{fig:linha}); para os anos
seguintes há relatos, que valem como indícios e não concordam entre si. Uma
entrevistada de \citeonline[p.~115]{aguiar2017} referiu-se à prática, em 2017,
como coisa de alguns anos antes, e uma agricultora relatou em 2016 que não
houve aplicações a partir de 2012, pela falta de chuvas \cite{nogueira2016}. Em sentido contrário, outra entrevistada do mesmo
estudo relata muitos episódios entre 2010 e 2014
\cite[p.~108--109]{aguiar2017}, e o parecer técnico entregue à Assembleia em
2016 afirma, no presente, que a pulverização aérea ocorre em Limoeiro do Norte,
Quixeré e Russas \cite{tramas2016parecer}.

\begin{figure}[htb]
\centering
\caption{Linha do tempo da pulverização aérea na Chapada do Apodi cearense}
\label{fig:linha}
\begin{tikzpicture}[x=0.56cm, y=1cm, >=Stealth, every node/.style={font=\footnotesize}]
  % eixo: 2000 = 0 ... 2026 = 26
  \fill[gray!55] (0,0) rectangle (10,0.45);
  \fill[gray!22] (10,0) rectangle (19.02,0.45);
  \fill[white]   (19.02,0) rectangle (24.97,0.45);
  \fill[pattern=north east lines, pattern color=gray!70] (24.97,0) rectangle (26,0.45);
  \draw (0,0) rectangle (26,0.45);
  \draw (10,0) -- (10,0.45); \draw (19.02,0) -- (19.02,0.45); \draw (24.97,0) -- (24.97,0.45);
  \node at (5,0.225) {A}; \node at (14.5,0.225) {B};
  \node at (22,0.225) {C}; \node at (25.5,0.225) {D};
  \foreach \a/\x in {2000/0,2005/5,2010/10,2015/15,2020/20,2025/25}
    {\draw (\x,0) -- (\x,-0.12); \node[below] at (\x,-0.12) {\a};}
  \draw[->] (15.15,1.0) -- (15.15,0.5);
  \node[above, align=center] at (14.3,1.0) {projeto de lei\\(2015)};
  \draw[->] (19.02,1.0) -- (19.02,0.5);
  \node[above, align=center] at (19.6,1.0) {Lei 16.820\\(2019)};
  \draw[->] (24.97,1.0) -- (24.97,0.5);
  \node[above, align=center] at (24.6,1.0) {Lei 19.135\\(12/2024)};
\end{tikzpicture}
\fonte{Elaboração própria, com base nas fontes citadas nesta seção. A:
aplicação aérea documentada. B: relatos divergentes; registros a obter. C:
proibição integral. D: proibição com exceção para drones.}
\end{figure}


Medir o que havia é parte do projeto. A norma federal prevê relatório mensal
dos operadores ao Ministério da Agricultura \cite{olivete2025}, e o
receituário agronômico registra a modalidade prescrita. O projeto solicita
esses registros e fixa antes os critérios que separam ausência de voo de
ausência de registro (Apêndice~\ref{ap:plano} e Seção~\ref{sec:robustez}).
